\documentclass[a4paper,12pt]{article}
\usepackage[margin=1.5cm]{geometry}
\usepackage{amsmath, amssymb}
\usepackage{array}
\usepackage{alphalph}

\usepackage{amsmath, amssymb, amsfonts}
\usepackage{multicol}
\usepackage{multirow}

\usepackage{enumitem}
\usepackage{booktabs}
\usepackage{tabularx}
\usepackage{array}
\usepackage{framed}
\usepackage{xcolor}
\usepackage{tikz}
\usepackage{fancyhdr}

% Estilos para cajas de teoría y recordatorios
\definecolor{graybox}{gray}{0.92}
\newenvironment{boxsummary}{%
  \def\FrameCommand{\fcolorbox{black}{graybox}}%
  \MakeFramed{\FrameRestore}%
  \small\textbf{RECORDAR:}\par\smallskip
}{%
  \endMakeFramed
}

\begin{document}

\section*{Ejercicio 1: Factoriza los siguientes polinomios}

\subsection*{a}

     \[
    3x^2+2x-8
    \]

    \[
    x = \dfrac{-2 \pm \sqrt{2^2 - 4 \cdot 3 \cdot (-8)}}{2 \cdot 3} = \dfrac{-2 \pm \sqrt{100}}{6} -> x_1 = 4/3 y x_2 = -2 
    \]

    \[
    = (x+2) \cdot 3 \cdot (x-\dfrac{4}{3})=(x+2)(3x-4)
    \]


\subsection*{b}
Primero, sacamos factor común $3x$:

\begin{align*}
3x^5-48x
&=3x(x^4-16)
\end{align*}

Dentro del paréntesis tenemos una diferencia de cuadrados,
ya que $x^4=(x^2)^2$ y $16=4^2$.
Aplicamos la identidad:

\[
a^2-b^2=(a-b)(a+b)
\]

Así, obtenemos:

\begin{align*}
3x(x^4-16)
&=3x(x^2-4)(x^2+4)\\
&=3x(x-2)(x+2)(x^2+4)
\end{align*}

Por tanto, la factorización es:

\[
\boxed{
3x^5-48x=3x(x-2)(x+2)(x^2+4)
}
\]

\subsection*{c}

Buscamos una raíz entera. Probamos $x=2$:

\[
2(2)^3+(2)^2-5(2)-10=16+4-10-10=0
\]

Como el resultado es cero, $x-2$ es un factor del polinomio.

Dividimos $2x^3+x^2-5x-10$ entre $x-2$.
Mediante la regla de Ruffini obtenemos el cociente
$2x^2+5x+5$.

Por tanto:

\begin{align*}
2x^3+x^2-5x-10
&=(x-2)(2x^2+5x+5)
\end{align*}

El discriminante del segundo factor es:

\[
\Delta=b^2-4ac=5^2-4(2)(5)=25-40=-15
\]

Como el discriminante es negativo, el trinomio no se
puede descomponer en factores lineales reales.

La factorización es:

\[
\boxed{
2x^3+x^2-5x-10=(x-2)(2x^2+5x+5)
}
\]

\subsection*{d}
Buscamos una raíz entera. Probamos $x=-1$:

\[
(-1)^3-7(-1)^2+8(-1)+16
=-1-7-8+16=0
\]

Como el resultado es cero, $x+1$ es un factor.

Dividimos el polinomio entre $x+1$ mediante Ruffini
y obtenemos el cociente $x^2-8x+16$.

Por tanto:

\begin{align*}
x^3-7x^2+8x+16
&=(x+1)(x^2-8x+16)
\end{align*}

El trinomio $x^2-8x+16$ es un cuadrado perfecto,
porque $16=4^2$ y $-8x=-2\cdot x\cdot4$.

Aplicamos la identidad:

\[
a^2-2ab+b^2=(a-b)^2
\]

Así:

\begin{align*}
(x+1)(x^2-8x+16)
&=(x+1)(x-4)^2
\end{align*}

La factorización es:

\[
\boxed{
x^3-7x^2+8x+16=(x+1)(x-4)^2
}
\]

\subsection*{e}

Observamos que todos los términos contienen $x$.
Sacamos factor común:

\[
x^4+2x^3-23x^2-60x
=x(x^3+2x^2-23x-60)
\]

Ahora buscamos una raíz entera del polinomio cúbico.
Probamos $x=5$:

\[
5^3+2(5)^2-23(5)-60
=125+50-115-60=0
\]

Por tanto, $x-5$ es un factor del polinomio cúbico.

Dividiendo entre $x-5$ mediante Ruffini obtenemos
el cociente $x^2+7x+12$.

Entonces:

\begin{align*}
x(x^3+2x^2-23x-60)
&=x(x-5)(x^2+7x+12)
\end{align*}

Factorizamos el trinomio $x^2+7x+12$.
Buscamos dos números cuyo producto sea $12$
y cuya suma sea $7$. Estos números son $3$ y $4$.

Por tanto:

\begin{align*}
x(x-5)(x^2+7x+12)
&=x(x-5)(x+3)(x+4)
\end{align*}

La factorización es:

\[
\boxed{
x^4+2x^3-23x^2-60x=x(x-5)(x+3)(x+4)
}
\]

\subsection*{f}

Buscamos una raíz entera. Probamos $x=1$:

\[
9(1)^4-36(1)^3+26(1)^2+4(1)-3
=9-36+26+4-3=0
\]

Como el resultado es cero, $x-1$ es un factor.

Dividimos el polinomio entre $x-1$ mediante Ruffini
y obtenemos el cociente
$9x^3-27x^2-x+3$.

Así:

\begin{align*}
9x^4-36x^3+26x^2+4x-3
&=(x-1)(9x^3-27x^2-x+3)
\end{align*}

Ahora agrupamos los términos del polinomio cúbico:

\begin{align*}
9x^3-27x^2-x+3
&=9x^2(x-3)-1(x-3)\\
&=(x-3)(9x^2-1)
\end{align*}

El término $9x^2-1$ es una diferencia de cuadrados:

\[
9x^2-1=(3x)^2-1^2=(3x-1)(3x+1)
\]

Sustituimos y obtenemos:

\begin{align*}
9x^4-36x^3+26x^2+4x-3
&=(x-1)(x-3)(9x^2-1)\\
&=(x-1)(x-3)(3x-1)(3x+1)
\end{align*}

La factorización es:

\[
\boxed{
9x^4-36x^3+26x^2+4x-3
=(x-1)(x-3)(3x-1)(3x+1)
}
\]

\section*{Resumen de los resultados}

\begin{align*}
\text{a)}\quad
3x^2+2x-8
&=(x+2)(3x-4)\\[8pt]
\text{b)}\quad
3x^5-48x
&=3x(x-2)(x+2)(x^2+4)\\[8pt]
\text{c)}\quad
2x^3+x^2-5x-10
&=(x-2)(2x^2+5x+5)\\[8pt]
\text{d)}\quad
x^3-7x^2+8x+16
&=(x+1)(x-4)^2\\[8pt]
\text{e)}\quad
x^4+2x^3-23x^2-60x
&=x(x-5)(x+3)(x+4)\\[8pt]
\text{f)}\quad
9x^4-36x^3+26x^2+4x-3
&=(x-1)(x-3)(3x-1)(3x+1)
\end{align*}

\newpage 
\subsection*{Simplifica la fracción}

$\dfrac{x^3-5x^2+6x}{x^3-x^2-14x+24} = \dfrac{ x \cdot (x-3) \cdot (x-2)}{(x - 2)(x - 3)(x + 4)} = \dfrac{x}{x+4}$

\begin{itemize}
    \item $x^3-5x^2+6x = x \cdot (x^2-5x+6) = x \cdot (x-3) \cdot (x-2) $
    
    \[ x= \dfrac{-b \pm \sqrt{b^2 - 4ac}}{2a} = \dfrac{5 \pm \sqrt{(-5)^2-4 \cdot 1 \cdot 6}}{2 \cdot 1 } = \dfrac{5 \pm \sqrt{1}}{2} => x_1= 3 \hspace{0.5cm} x_2=2\]

    \item $x^3-x^2-14x+24$
    
    El término independiente es $24$. Los posibles divisores enteros de $24$ ($\pm a$) son:
    \[
    \text{Divisores}(24) = \{\pm 1, \pm 2, \pm 3, \pm 4, \pm 6, \pm 8, \pm 12, \pm 24\}
    \]

    Encontrar la primera raíz utilizando el el teorea del resto en lugar de ruffini. 

    Evaluamos el polinomio $P(x)$ en los divisores hasta encontrar un valor que dé cero:

    \begin{itemize}
        \item Para $x = 1$:
        \[ P(1) = (1)^3 - (1)^2 - 14(1) + 24 = 1 - 1 - 14 + 24 = 10 \neq 0 \]
        \item Para $x = 2$:
        \[ P(2) = (2)^3 - (2)^2 - 14(2) + 24 = 8 - 4 - 28 + 24 = 0 \]
    \end{itemize}

    Como $P(2) = 0$, $x = 2$ es una raíz del polinomio, lo que significa que \textbf{$(x - 2)$} es un factor de $P(x)$.

    Dividimos $P(x)$ entre $(x - 2)$ usando la Regla de Ruffini con los coeficientes $[1, -1, -14, 24]$ y la raíz $x = 2$:

    \begin{center}
    \begin{tabular}{r|rrrr}
    & 1 & -1 & -14 & 24 \\
    2 &   &  2 &   2 & -24 \\
    \hline
    & 1 &  1 & -12 & \textbf{0}
    \end{tabular}
    \end{center}

    El cociente resultante es un polinomio de segundo grado:
    \[
    Q(x) = x^2 + x - 12
    \]
    Por lo tanto, la factorización parcial es:
    \[
    P(x) = (x - 2)(x^2 + x - 12)
    \]

    Ahora factorizamos $Q(x) = x^2 + x - 12$. 
    \[
    x = \dfrac{-1 \pm \sqrt{ 1 - 4 \cdot (-12)}}{2} = \dfrac{-1 \pm \sqrt{49}}{2} => x_1 = -4 \hspace{0.5cm}x_2=3
    \]

    Uniendo todos los factores obtenidos, la factorización completa del polinomio es:

    \[
    \mathbf{x^3 - x^2 - 14x + 24 = (x - 2)(x - 3)(x + 4)}
    \]
\end{itemize}

 








\end{document}


